\subsection{Curvature and projective complement exhaustions}
\label{subsec:projective-complement-curvature-index}

\Cref{cor:dual-norm-complement-convexity-obstruction} excludes $r$-convexity of the complement of the zero section in the projective completion, but it does not decide which larger indices occur. Upper bounds come from curvature. The reciprocal of the squared norm of a Hermitian metric is an exhaustion near the removed zero section, and every positive direction of the curvature, evaluated at a fiber vector, contributes a positive Levi eigenvalue. For a stable bundle of positive Gauduchon degree, the Hermitian--Einstein metric provides such a direction at every fiber vector.

In this subsection, $E\to X$ is a holomorphic vector bundle of rank $r\geq1$, $n\geq1$, and $h$ is a smooth Hermitian metric on $E$. In a holomorphic frame, write the matrix of $h$ as $M$, using column vectors and an inner product linear in its first argument, so that $h(v,v)=v^*Mv$. Put
\[
 F_h=\bar\partial(M^{-1}\partial M)
     =\sum_{\alpha,\beta}F_{\alpha\bar\beta}\,\mathrm dz_\alpha\wedge\mathrm d\bar z_\beta,
 \qquad R_h(u,\bar u)=\sum_{\alpha,\beta}F_{\alpha\bar\beta}u_\alpha\bar u_\beta.
\]
For $v\in E_x$, the real $(1,1)$-form $\alpha_{h,v}=\mathrm i\,h(F_hv,v)$ has the Hermitian coefficient form $u\mapsto h(R_h(u,\bar u)v,v)$. \index{Chern curvature!of a bundle metric}For a Hermitian form $\omega$ on $X$, the \emph{mean curvature} of $h$ is the $h$-Hermitian endomorphism $K_{\omega,h}=\Lambda_\omega(\mathrm iF_h)$, where $\Lambda_\omega$ denotes contraction with $\omega$, so that $\Lambda_\omega(\beta)\,\omega^n=n\beta\wedge\omega^{n-1}$ for $(1,1)$-forms $\beta$.\index{Mean curvature!of a bundle metric} The curvature $F_h$ is unscaled: for a line bundle, the curvature $-\ddc\log h$ of \cref{sec:convex-bundle-norm-obstruction} is $\mathrm iF_h/(2\pi)$.

Use the line convention for projectivization and put
\[
 U_E=\mathbb P(\mathcal O_X\oplus E)\setminus\mathbb P(\mathcal O_X).
\]
The chart $[1:v]$ identifies the total space of $E$ with an open subset of $\mathbb P(\mathcal O_X\oplus E)$, in which $\mathbb P(\mathcal O_X)$ is the zero section; thus $U_E$ contains $E^\times$. Its complex dimension is $n+r$. Convexity is understood in the sense of the smooth exhaustions of \cref{subsec:concave-neighbourhood-dual-norms}.

\Needspace{14\baselineskip}
\begin{lemma}\label{lem:reciprocal-bundle-norm-levi}
\index{Chern curvature!reciprocal norm Hessian}
Put $G(v)=h(v,v)$ and $\rho=G^{-1}$ on $E^\times$. In a holomorphic frame that is unitary and \emph{Chern-normal} at $x$, meaning that $M(x)$ is the identity and the first derivatives of $M$ vanish at $x$, the complex Hessian of $\rho$ at $(x,v)$ is
\[
 \mathcal L_\rho(u,w)
 =\frac{h(R_h(u,\bar u)v,v)}{G^2}
  -\frac{|w|_h^2}{G^2}
  +\frac{2|h(w,v)|^2}{G^3},
 \qquad u\in T_x^{1,0}X,\quad w\in E_x,
\]
where $(u,w)$ are the horizontal and vertical components of a tangent vector. There are no mixed horizontal--vertical terms, and the vertical block has exactly one positive and $r-1$ negative eigenvalues.
\end{lemma}
\begin{proof}
Such a frame is obtained by a constant change of frame making it unitary at $x$, followed by a holomorphic change of frame, linear in the coordinates, that cancels $\partial M$. Then $F_{\alpha\bar\beta}=-\partial_\alpha\partial_{\bar\beta}M$ at $x$. The horizontal first derivatives of $G$ and its mixed horizontal--vertical second derivatives vanish. The horizontal complex Hessian of $G$ is $-h(R_h(u,\bar u)v,v)$, its vertical complex Hessian is $|w|_h^2$, and its vertical first derivative in the direction $w$ is $h(w,v)$. Apply
\[
 \partial\bar\partial(G^{-1})
 =2G^{-3}\partial G\wedge\bar\partial G-G^{-2}\partial\bar\partial G.
\]
The vertical block equals $G^{-1}>0$ at $w=v$ and equals $-|w|_h^2/G^2$ when $w\perp_h v$. These two subspaces are orthogonal for the vertical block, which proves the statement about its eigenvalues.
\end{proof}

\Needspace{8\baselineskip}
\begin{theorem}\label{thm:bundle-curvature-complement-convexity}
\index{Projective bundle!curvature exhaustion}\index{$q$-convexity!bundle complements}
Let $s$ be an integer with $0\leq s\leq n$. Suppose that for every $x\in X$ and every $0\ne v\in E_x$, the coefficient form of $\alpha_{h,v}$ has at least $s$ positive eigenvalues. Then $U_E$ is $(n+r-s)$-convex.
\end{theorem}
\begin{proof}
By \cref{lem:reciprocal-bundle-norm-levi}, the horizontal block of $\mathcal L_\rho$ has at least $s$ positive eigenvalues and the vertical block has one. Their direct sum thus has at least $s+1$ positive eigenvalues at every point of $E^\times$.

Choose $\epsilon>0$ and a smooth function $\chi:[0,+\infty)\to[0,1]$ equal to one on $[0,\epsilon]$ and to zero on $[2\epsilon,+\infty)$. Define
\[
 \psi([1:v])=\frac{\chi(G(v))}{G(v)}\quad(v\ne0),
\]
and extend $\psi$ by zero on the divisor at infinity $\mathbb P(E)$. This extension is smooth. Indeed, near a point $[0:v]$ with $v\ne0$, use a local representative $[u:v]$ in a projective chart. On its intersection with the affine chart, $G(v/u)=h(v,v)/|u|^2$ tends to infinity as $u\to0$, so $\chi(G)$ vanishes in a neighborhood of the point.

The function $\psi$ tends to infinity near the removed zero section. Since the ambient projective bundle is compact, every sublevel set of $\psi$ is compact in $U_E$: it is closed there and cannot accumulate on the removed section. Put $D=\{G<\epsilon\}\subset E$, including the zero section, and $K=\mathbb P(\mathcal O_X\oplus E)\setminus D$. Then $K$ is a compact subset of $U_E$, and $\psi=\rho$ on $U_E\setminus K$. Thus the exhaustion $\psi$ has at least $s+1$ positive Levi eigenvalues outside a compact set. In dimension $n+r$, this is $(n+r-s)$-convexity.
\end{proof}

Taking $s=0$ gives $(n+r)$-convexity for every bundle. For $s=1$, the curvature hypothesis is precisely RC-positivity of the metric, as in \cref{def:directional-rc-positivity}. A positive definite mean curvature provides such a direction by taking a trace.

To apply this to stable bundles, we extend the degree of line bundles of \cref{sec:convex-bundle-norm-obstruction} to vector bundles through their determinants.

\Needspace{13\baselineskip}
\begin{definition}\label{def:gauduchon-bundle-slope}
\index{Gauduchon degree!vector bundles}\index{Slope stability!Gauduchon form}\index{Bott--Chern class!Gauduchon degree}
For a Gauduchon form $\omega$ on $X$, define
\[
 \deg_\omega E=\deg_{\omega^{n-1}}\det E
 =\int_X\frac{\mathrm i\operatorname{tr}F_h}{2\pi}\wedge\omega^{n-1},
 \qquad \mu_\omega(E)=\frac{\deg_\omega E}{r}.
\]
For a nonzero torsion-free coherent sheaf $\mathcal F$, use the determinant line bundle $\det\mathcal F=(\bigwedge^{\operatorname{rank}\mathcal F}\mathcal F)^{**}$ to define its degree and slope. The bundle $E$ is \emph{$\omega$-slope-stable} if $\mu_\omega(\mathcal F)<\mu_\omega(E)$ for every coherent subsheaf $\mathcal F\subset E$ with $0<\operatorname{rank}\mathcal F<r$ and torsion-free quotient.
\end{definition}

If $\det h'=\mathrm e^f\det h$, the two curvature representatives differ by $-\ddc f$, whose integral against $\omega^{n-1}$ vanishes by \cref{lem:ibp}. Thus the degree depends only on the Bott--Chern class $c_1^{\mathrm{BC}}(E)$ of the closed form $\mathrm i\operatorname{tr}F_h/(2\pi)$, and it is the pairing of this class with $\omega^{n-1}$. The Hopf example in \cref{ex:hopf-exact-complement-index} shows that the de Rham first Chern class of $E$ does not determine the degree.

\Needspace{10\baselineskip}
\begin{corollary}\label{cor:positive-mean-complement-convexity}
\index{Hermitian--Einstein metric!positive degree}\index{Mean curvature!complement convexity}
If $K_{\omega,h}$ is positive definite for a Hermitian form $\omega$ on $X$, then $U_E$ is $(n+r-1)$-convex. This holds, in particular, when $\omega$ is Gauduchon, $E$ is $\omega$-slope-stable, and $\deg_\omega E>0$.
\end{corollary}
\begin{proof}
For $v\ne0$, the trace of $\alpha_{h,v}$ with respect to $\omega$ is $h(K_{\omega,h}v,v)>0$, and a Hermitian form with positive trace has a positive eigenvalue. Apply \cref{thm:bundle-curvature-complement-convexity} with $s=1$.

Under the Gauduchon and stability hypotheses, the classical existence theorem for Hermitian--Einstein metrics on stable bundles gives a metric $h$ with $K_{\omega,h}=\lambda\operatorname{Id}_E$, where $\lambda$ is a constant. Taking the trace and integrating the contraction identity gives
\[
 r\lambda\int_X\omega^n
 =n\int_X\operatorname{tr}(\mathrm iF_h)\wedge\omega^{n-1}
 =2\pi n\deg_\omega E>0.
\]
Here $\operatorname{tr}F_h=\bar\partial\partial\log\det M$, so $\mathrm i\operatorname{tr}F_h/(2\pi)$ represents $c_1^{\mathrm{BC}}(E)$. Thus $\lambda>0$, and the first assertion applies.
\end{proof}

\Needspace{9\baselineskip}
\begin{corollary}\label{cor:hirzebruch-exact-complement-index}
\index{Hirzebruch surface!least complement convexity index}
For every $e\geq0$, there is a stable ample rank-two bundle $E$ on $X=\mathbb F_e$ as in \cref{cor:hirzebruch-non-two-convex-complement} for which $U_E$ is $3$-convex but not $2$-convex. Its least smooth convexity index is three.
\end{corollary}
\begin{proof}
Use the ruling section $S$ and fiber $f$ of \cref{lem:hirzebruch-ruling-sections}, and put $H=S+(3e+1)f$ and $P=4H-f$. Choose $E$ as in \cref{cor:hirzebruch-non-two-convex-complement}, $\mu_P$-stable and with $c_1(E)=4H$. The intersection numbers of \cref{lem:hirzebruch-ruling-sections} give
\[
 c_1(E)\cdot P=4H\cdot(4H-f)=80e+28>0.
\]
By the same lemma, $P=4S+(12e+3)f$ is very ample. Take for $\omega$ the pulled-back Fubini--Study K\"ahler form representing $c_1(\mathcal O_X(P))$. Degrees with respect to $\omega$ are intersection numbers with $P$, so $E$ is $\omega$-slope-stable with $\deg_\omega E>0$, and \cref{cor:positive-mean-complement-convexity} gives $3$-convexity, since $n=r=2$. By the choice of $E$, the complement is not $2$-convex. Finally, $1$-convexity would imply $2$-convexity, so the least index is three.
\end{proof}
